% Generated by paper/make_numbers.py from results/*.json.
% Do not edit by hand: the next run overwrites this file.
\begin{table}[t]
\centering
\caption{Held-out accuracy of networks retrained with the first $k$ blocks, by divergence level and normalization. The columns are receptive fields in positions. Means over \nSeeds\ seeds; this is the grid behind Figure~\ref{fig:reach}a and Table~\ref{tab:difficulty}.}
\label{tab:grid}
\resizebox{\ifdim\width>\linewidth\linewidth\else\width\fi}{!}{%
\begin{tabular}{llrrrrr}
\toprule
normalization & $F_{ST}$ & $R=9$ & $R=33$ & $R=129$ & $R=513$ & $R=2049$ \\
\midrule
pooled over length & 0.2405 & 0.9631 & 0.9817 & 0.9962 & 0.9982 & 0.9983 \\
 & 0.0719 & 0.9353 & 0.9478 & 0.9686 & 0.9836 & 0.9877 \\
 & 0.0380 & 0.9158 & 0.9250 & 0.9381 & 0.9532 & 0.9610 \\
 & 0.0199 & 0.8502 & 0.8649 & 0.8745 & 0.8907 & 0.8969 \\
 & 0.0096 & 0.7272 & 0.7339 & 0.7399 & 0.7507 & 0.7592 \\
per position & 0.2405 & 0.7228 & 0.8786 & 0.9849 & 0.9980 & 0.9983 \\
 & 0.0719 & 0.6204 & 0.7187 & 0.8618 & 0.9674 & 0.9862 \\
 & 0.0380 & 0.5767 & 0.6463 & 0.7644 & 0.8914 & 0.9508 \\
 & 0.0199 & 0.5521 & 0.5977 & 0.6765 & 0.7763 & 0.8720 \\
 & 0.0096 & 0.5190 & 0.5487 & 0.5945 & 0.6509 & 0.7200 \\
\bottomrule
\end{tabular}}
\end{table}

\begin{table}[t]
\centering
\caption{Reach worth at each divergence level of the simulated genomes, with statistics pooled along the sequence and per position. Their ratio, per position over pooled, does not vary systematically with difficulty.}
\label{tab:difficulty}
\resizebox{\ifdim\width>\linewidth\linewidth\else\width\fi}{!}{%
\begin{tabular}{lrrr}
\toprule
$F_{ST}$ & pooled over length & per position & ratio \\
\midrule
0.2405 & $+0.0351 \pm 0.0019$ & $+0.2754 \pm 0.0099$ & 7.8 \\
0.0719 & $+0.0524 \pm 0.0110$ & $+0.3658 \pm 0.0051$ & 7.0 \\
0.0380 & $+0.0452 \pm 0.0107$ & $+0.3741 \pm 0.0052$ & 8.3 \\
0.0199 & $+0.0467 \pm 0.0077$ & $+0.3199 \pm 0.0126$ & 6.8 \\
0.0096 & $+0.0320 \pm 0.0119$ & $+0.2011 \pm 0.0173$ & 6.3 \\
\bottomrule
\end{tabular}}
\end{table}

\begin{table}[t]
\centering
\caption{Switch-rate dose--response on both simulated processes and on real haplotypes. Both middle columns are the accuracy gained by the \reachFold-fold enlargement of the receptive field; the last column is the constant of the $c/n$ regularity at each density. It is flat within each process and differs between them.}
\label{tab:dose}
\resizebox{\ifdim\width>\linewidth\linewidth\else\width\fi}{!}{%
\begin{tabular}{llccrr}
\toprule
process & $n$ & pooled over length & per position & ratio & $(\text{ratio}-1)\,n$ \\
\midrule
genomic & 0.29 & $+0.0452 \pm 0.0107$ & $+0.3741 \pm 0.0052$ & 8.3 & 2.08 \\
 & 0.99 & $+0.1098 \pm 0.0070$ & $+0.3453 \pm 0.0074$ & 3.1 & 2.13 \\
 & 3.00 & $+0.1745 \pm 0.0026$ & $+0.2958 \pm 0.0061$ & 1.7 & 2.08 \\
 & 9.89 & $+0.1692 \pm 0.0089$ & $+0.2151 \pm 0.0062$ & 1.3 & 2.68 \\
synthetic & 0.09 & $+0.0062 \pm 0.0014$ & $+0.3892 \pm 0.0038$ & 62.8 & 5.80 \\
 & 0.29 & $+0.0246 \pm 0.0040$ & $+0.3740 \pm 0.0047$ & 15.2 & 4.05 \\
 & 0.99 & $+0.0785 \pm 0.0078$ & $+0.3349 \pm 0.0034$ & 4.3 & 3.25 \\
 & 2.98 & $+0.1245 \pm 0.0039$ & $+0.2630 \pm 0.0083$ & 2.1 & 3.32 \\
 & 5.99 & $+0.1269 \pm 0.0071$ & $+0.1990 \pm 0.0036$ & 1.6 & 3.40 \\
 & 9.89 & $+0.1094 \pm 0.0008$ & $+0.1543 \pm 0.0043$ & 1.4 & 4.05 \\
 & 19.87 & $+0.0860 \pm 0.0041$ & $+0.1022 \pm 0.0003$ & 1.2 & 3.74 \\
real 1000G & 0.13 & $+0.0181 \pm 0.0105$ & $+0.2190 \pm 0.0195$ & 12.1 & 1.40 \\
\bottomrule
\end{tabular}}
\end{table}

\begin{table}[t]
\centering
\caption{The Conv-TasNet separator on the synthetic process at 0.27 switches per sequence, against an exact whole-sequence optimum of 0.9522. Its depthwise kernel is 3, so a block adds $2d$ rather than $4d$ and the receptive fields differ from the other tables. Means over \nSeeds\ seeds.}
\label{tab:tasnet}
\resizebox{\ifdim\width>\linewidth\linewidth\else\width\fi}{!}{%
\begin{tabular}{llrrrrr}
\toprule
layer & statistics pooled over & $R=7$ & $R=19$ & $R=67$ & $R=259$ & $R=1027$ \\
\midrule
gLN & channels and time & 0.9246 & 0.9244 & 0.9245 & 0.9289 & 0.9407 \\
cLN & channels and time up to $t$ & 0.8676 & 0.8705 & 0.8740 & 0.8860 & 0.9186 \\
per position & channels & 0.5423 & 0.5681 & 0.6262 & 0.7314 & 0.8745 \\
\bottomrule
\end{tabular}}
\end{table}

\begin{table}[t]
\centering
\caption{(P4) A transformer, whose attention spans the input sequence at every depth, with sequence-pooled and with per-position normalization. The two are within a few thousandths of each other and of the exact optimum at every switch density: with the context already available there is nothing for the path to supply. The last column is the better of the two normalizations against the optimum, and the range the main text quotes is the widest of these. Means over \nSeeds\ seeds.}
\label{tab:transformer}
\resizebox{\ifdim\width>\linewidth\linewidth\else\width\fi}{!}{%
\begin{tabular}{rrrrrr}
\toprule
switches per sequence & pooled over length & per position & difference & exact optimum & shortfall \\
\midrule
0.27 & 0.9497 & 0.9472 & +0.0025 & 0.9522 & 0.0025 \\
0.98 & 0.8956 & 0.8932 & +0.0024 & 0.8993 & 0.0037 \\
2.98 & 0.8146 & 0.8125 & +0.0021 & 0.8194 & 0.0047 \\
9.79 & 0.7023 & 0.7000 & +0.0022 & 0.7058 & 0.0035 \\
\bottomrule
\end{tabular}}
\end{table}

\begin{table}[t]
\centering
\caption{Per-position networks against the local bound for their own receptive field, on the synthetic process: each entry is held-out accuracy minus that bound, so zero is the bound and positive is above it. This is the grid behind every number quoted in Section~\ref{sec:tests} for this comparison: \boundAllCells\ cells, mean absolute value \boundAllMean, \boundAllAbove\ of them positive, the largest \boundAllAboveMax. Exceedances are expected, since $\mathrm{local}(R)$ is a blocked bound and therefore conservative, and the entries are finite-sample means (Section~\ref{sec:setup}). Adding the corresponding entry of Table~\ref{tab:optima} recovers the accuracy itself. Means over \nSeeds\ seeds.}
\label{tab:resid}
\resizebox{\ifdim\width>\linewidth\linewidth\else\width\fi}{!}{%
\begin{tabular}{rrrrrr}
\toprule
switches/sequence & $R=9$ & $R=33$ & $R=129$ & $R=513$ & $R=2049$ \\
\midrule
0.09 & -0.0006 & -0.0017 & -0.0050 & -0.0124 & -0.0172 \\
0.29 & -0.0003 & -0.0016 & -0.0049 & -0.0115 & -0.0119 \\
0.99 & -0.0008 & -0.0016 & -0.0030 & -0.0049 & -0.0013 \\
2.98 & -0.0010 & -0.0007 & +0.0012 & +0.0019 & -0.0001 \\
5.99 & -0.0008 & +0.0005 & +0.0031 & +0.0084 & -0.0009 \\
9.89 & -0.0003 & +0.0010 & +0.0051 & +0.0091 & -0.0002 \\
19.87 & -0.0005 & +0.0011 & +0.0055 & +0.0045 & -0.0005 \\
\bottomrule
\end{tabular}}
\end{table}

\begin{table}[t]
\centering
\caption{The reference values on the synthetic process: the local bound at every receptive field it was computed for, and the exactly computed optima with full context and with an oracle given only the sequence's class proportion. The intermediate columns are here because the main text reads values off them: $R=33$ brackets the U-Net's receptive field (Section~\ref{sec:tests}) and $R=129$ is where the per-position control matches a sequence-pooled model at $R=9$ (Section~\ref{sec:tests}).}
\label{tab:optima}
\resizebox{\ifdim\width>\linewidth\linewidth\else\width\fi}{!}{%
\begin{tabular}{rrrrrrrr}
\toprule
switches & $R=9$ & $R=33$ & $R=129$ & $R=513$ & $R=2049$ & full & proportion oracle \\
\midrule
0.09 & 0.5479 & 0.5908 & 0.6741 & 0.8165 & 0.9537 & 0.9759 & 0.9760 \\
0.29 & 0.5479 & 0.5905 & 0.6738 & 0.8110 & 0.9335 & 0.9504 & 0.9348 \\
0.99 & 0.5482 & 0.5903 & 0.6708 & 0.7951 & 0.8835 & 0.8986 & 0.8390 \\
2.98 & 0.5493 & 0.5913 & 0.6660 & 0.7666 & 0.8115 & 0.8194 & 0.7239 \\
5.99 & 0.5474 & 0.5875 & 0.6543 & 0.7211 & 0.7466 & 0.7508 & 0.6567 \\
9.89 & 0.5466 & 0.5846 & 0.6424 & 0.6869 & 0.7008 & 0.7031 & 0.6290 \\
19.87 & 0.5470 & 0.5821 & 0.6244 & 0.6439 & 0.6492 & 0.6508 & 0.5858 \\
\bottomrule
\end{tabular}}
\end{table}

\begin{table}[t]
\centering
\caption{The simulated ancestry process against the benchmark it was ported from \citep{tian2026lai}, at all eight of its split times. $F_{ST}$ is measured from the simulated panels rather than set, so agreement is a check on the demography, the coalescent draw and the estimator together. Accuracy is the reference architecture's, over \nSeeds\ seeds.}
\label{tab:port}
\resizebox{\ifdim\width>\linewidth\linewidth\else\width\fi}{!}{%
\begin{tabular}{rrrrr}
\toprule
split time & $F_{ST}$, ours & $F_{ST}$, benchmark & difference & accuracy \\
\midrule
3200 & 0.2405 & 0.2434 & -0.0029 & 0.9983 \\
1600 & 0.1366 & 0.1376 & -0.0010 & 0.9955 \\
800 & 0.0719 & 0.0736 & -0.0017 & 0.9877 \\
400 & 0.0380 & 0.0384 & -0.0004 & 0.9610 \\
200 & 0.0199 & 0.0200 & -0.0001 & 0.8969 \\
100 & 0.0096 & 0.0097 & -0.0001 & 0.7592 \\
50 & 0.0053 & 0.0050 & +0.0003 & 0.6579 \\
25 & 0.0026 & 0.0030 & -0.0004 & 0.5721 \\
\bottomrule
\end{tabular}}
\end{table}

\begin{table}[t]
\centering
\caption{The U-Net depth sweep on the synthetic process, at \unetDepths\ depths and two switch densities, means over \nSeeds\ seeds. ``Receptive field'' is the convolutional span, measured by autograd rather than derived; with statistics pooled along the sequence the same measurement returns \unetWindow\ at every depth, which is the whole sequence. Depth is what sets the receptive field here, and it also sets how many values each normalization statistic is taken over, so this is the one architecture in which the two cannot be varied separately.}
\label{tab:unet}
\resizebox{\ifdim\width>\linewidth\linewidth\else\width\fi}{!}{%
\begin{tabular}{rrrrr}
\toprule
switches/sequence & depth & receptive field & pooled over length & per position \\
\midrule
0.27 & 1 & 22 & 0.9247 & 0.5730 \\
 & 2 & 48 & 0.9252 & 0.6060 \\
 & 3 & 96 & 0.9255 & 0.6493 \\
 & 4 & 192 & 0.9276 & 0.7039 \\
 & 5 & 384 & 0.9326 & 0.7756 \\
 & 6 & 768 & 0.9417 & 0.8482 \\
2.98 & 1 & 22 & 0.6964 & 0.5746 \\
 & 2 & 48 & 0.7066 & 0.6072 \\
 & 3 & 96 & 0.7242 & 0.6492 \\
 & 4 & 192 & 0.7504 & 0.6985 \\
 & 5 & 384 & 0.7808 & 0.7512 \\
 & 6 & 768 & 0.8050 & 0.7928 \\
\bottomrule
\end{tabular}}
\end{table}

\begin{sidewaystable}
\centering
\caption{Normalization configurations of 13 published sequence-labeling models across 4 domains whose labels come in long runs, read from each project's source rather than from its paper. ``Exposed'' means the criterion of Section~\ref{sec:theory} is met: statistics computed from the current input, along the sequence, at inference. Two are. Citations are to each model's publication; the evidence for the row is the code, since the two can disagree (Section~\ref{sec:discussion}). Each row's source link, pinned to the commit that was read, and its provenance are in the released \texttt{exposure\_survey.json}.}
\label{tab:survey}
\resizebox{\ifdim\width>\linewidth\linewidth\else\width\fi}{!}{%
\begin{tabular}{lllc}
\toprule
domain & model & normalization & exposed \\
\midrule
action segmentation & MS-TCN \citep{farha2019mstcn} & none & no \\
action segmentation & ASFormer \citep{yi2021asformer} & InstanceNorm1d, no running stats & \textbf{yes} \\
sleep staging & U-Sleep \citep{perslev2021usleep} & BatchNorm, running stats at eval & no \\
sleep staging & U-Time \citep{perslev2019utime} & BatchNorm, running stats at eval & no \\
sleep staging & DeepSleepNet \citep{supratak2017deepsleepnet} & BatchNorm, moving stats at eval & no \\
sleep staging & TinySleepNet \citep{supratak2020tinysleepnet} & BatchNorm, frozen at eval & no \\
sleep staging & SleepTransformer \citep{phan2022sleeptransformer} & LayerNorm, feature axis & no \\
speaker diarization & SA-EEND \citep{fujita2019eend} & LayerNorm, feature axis & no \\
speaker diarization & EEND-EDA \citep{horiguchi2020eendeda} & LayerNorm, feature axis & no \\
speaker diarization & ECAPA-TDNN \citep{desplanques2020ecapa} & BatchNorm, running stats at eval & no \\
speaker diarization & pyannote.audio SincNet frontend \citep{bredin2020pyannote} & InstanceNorm1d, no running stats & \textbf{yes} \\
genomics & Enformer \citep{avsec2021effective} & BatchNorm, moving averages; LayerNorm, feature axis & no \\
genomics & Basenji \citep{kelley2018sequential} & BatchNorm, moving averages; LayerNorm, feature axis & no \\
\bottomrule
\end{tabular}}
\end{sidewaystable}

\begin{table}[t]
\centering
\caption{Accuracy lost by removing the long-range blocks (those that enlarge the receptive field; Section~\ref{sec:setup}), per divergence level and normalization: ablating them from the trained network against retraining without them. The pooled figures of Section~\ref{sec:ablation} are means over these rows: \ablCostOpen\ against \retCostOpen\ (ratio \ablRatioOpen) with statistics pooled along the sequence, and a ratio of \ablRatioClosed\ per position. Means over \nSeeds\ seeds.}
\label{tab:ablcost}
\resizebox{\ifdim\width>\linewidth\linewidth\else\width\fi}{!}{%
\begin{tabular}{llrrr}
\toprule
normalization & $F_{ST}$ & ablated & retrained & ratio \\
\midrule
pooled over length & 0.2405 & 0.3117 & 0.0351 & 8.9 \\
 & 0.0719 & 0.4273 & 0.0524 & 8.2 \\
 & 0.0380 & 0.4099 & 0.0452 & 9.1 \\
 & 0.0199 & 0.3774 & 0.0467 & 8.1 \\
 & 0.0096 & 0.2289 & 0.0320 & 7.2 \\
per position & 0.2405 & 0.4085 & 0.2754 & 1.5 \\
 & 0.0719 & 0.4491 & 0.3658 & 1.2 \\
 & 0.0380 & 0.4371 & 0.3741 & 1.2 \\
 & 0.0199 & 0.3584 & 0.3199 & 1.1 \\
 & 0.0096 & 0.2084 & 0.2011 & 1.0 \\
\bottomrule
\end{tabular}}
\end{table}

\begin{table}[t]
\centering
\caption{The derived advantage of an oracle given the sequence's class proportion, $1/\sqrt{2\pi n}$, against the value computed exactly. No quantity is fitted. The approximation requires $n \gg 1$ and $p \ll 1$ (Section~\ref{sec:theory}) and fails below $n \approx 1$ because $\bar\pi$ is bounded in $[0,1]$.}
\label{tab:deriv}
\begin{tabular}{rrrr}
\toprule
switches per sequence & computed & $1/\sqrt{2\pi n}$ & error \\
\midrule
0.09 & 0.4760 & 1.3029 & -63\% \\
0.29 & 0.4348 & 0.7471 & -42\% \\
0.99 & 0.3390 & 0.4001 & -15\% \\
2.98 & 0.2239 & 0.2311 & -3\% \\
5.99 & 0.1567 & 0.1631 & -4\% \\
9.89 & 0.1290 & 0.1269 & +2\% \\
19.87 & 0.0858 & 0.0895 & -4\% \\
\bottomrule
\end{tabular}
\end{table}
